% Desarrollo con asistencia de IA; pendiente de revisión del integrante responsable.
\subsection{Tiempo físico, causalidad y concurrencia}
La hora del sistema permite fechar un registro, pero no basta para determinar
qué evento causó otro: los relojes físicos pueden diferir y el orden de llegada
depende de la red y de la planificación de procesos. Los relojes lógicos
representan restricciones de orden, no segundos transcurridos
\cite{vanSteenTanenbaum2023}.

La relación de precedencia causal, escrita $a\rightarrow b$, se define por:
\begin{enumerate}
  \item El orden de los eventos dentro de un mismo proceso.
  \item El envío de un mensaje antes de su recepción.
  \item La transitividad: si $a\rightarrow b$ y $b\rightarrow c$, entonces
  $a\rightarrow c$.
\end{enumerate}
Dos eventos son concurrentes cuando no se cumple ninguna de las relaciones
$a\rightarrow b$ ni $b\rightarrow a$. Esta concurrencia es causal: no exige
que ambos eventos ocurran en el mismo instante físico.

\subsection{Reglas del reloj de Lamport}
Cada nodo $i$ mantiene un contador $L_i$, inicialmente cero. Antes de registrar
un evento local o enviar un mensaje, realiza:
\begin{equation}
L_i \leftarrow L_i+1.
\label{eq:lamport-envio}
\end{equation}
El mensaje transporta el contador resultante. Si se recibe una marca $t$, la
actualización es:
\begin{equation}
L_i \leftarrow \max(L_i,t)+1.
\label{eq:lamport-recepcion}
\end{equation}
La recepción y el envío de su acuse son eventos distintos; por ello, el acuse
requiere otro incremento. Estas reglas garantizan la condición del reloj:
\begin{equation}
a\rightarrow b \;\Longrightarrow\; L(a)<L(b).
\end{equation}
El recíproco es falso: $L(a)<L(b)$ no demuestra que exista una cadena causal
entre ambos. Tampoco una diferencia de cinco unidades representa cinco
mensajes ni una duración física de cinco unidades de tiempo
\cite{vanSteenTanenbaum2023}.

Dos eventos de nodos distintos con la misma marca no pueden estar relacionados
por precedencia causal si se cumplen las reglas anteriores. Sin embargo, una
marca escalar no permite decidir la concurrencia para todos los pares de eventos:
también puede haber eventos concurrentes con marcas diferentes. Ordenar por
la pareja $(L,\mathrm{id\ del\ nodo})$ permite desempatar de forma determinista,
pero ese orden total no convierte los eventos en causalmente relacionados.

\subsection{Comparación con relojes vectoriales}
En un sistema de tres nodos, un reloj vectorial mantiene tres componentes,
una por nodo. En un evento local o envío se incrementa la componente propia;
al recibir se combina el vector local y el recibido mediante máximos por
componente y después se incrementa la componente del receptor.
Para dos vectores $V(a)$ y $V(b)$, se escribe $V(a)<V(b)$ si todas las
componentes del primero son menores o iguales y al menos una es estrictamente
menor. Bajo el modelo habitual de procesos y mensajes, esta comparación
representa la precedencia causal; vectores incomparables identifican
concurrencia \cite{vanSteenTanenbaum2023}.

El costo es mantener y transmitir información proporcional al número de nodos,
frente a un contador escalar en Lamport. Los relojes vectoriales se presentan
como comparación conceptual; no forman parte del módulo implementado.

\subsection{Ejemplo didáctico, no experimental}
Supóngase que A y C, con contadores inicialmente en cero, envían sus primeros
mensajes a B sin haber intercambiado información entre sí. Ambos envíos llevan
marca 1 y son concurrentes. Si B recibe primero el mensaje de A, actualiza su
contador a 2; al enviar el acuse lo incrementa a 3. A recibe ese acuse y pasa
a $\max(1,3)+1=4$. Si luego B recibe el mensaje inicial de C, con marca 1,
su contador pasa a $\max(3,1)+1=4$.

En este ejemplo, el envío de A precede causalmente a su recepción en B,
con $1<2$. Los primeros envíos de A y C tienen la misma marca sin relación
causal entre ellos. Los números anteriores ilustran las reglas; no son
mediciones ni sustituyen la comprobación sobre los CSV de la práctica.
